\documentclass[a4paper,11pt]{article}
\usepackage[T2A]{fontenc}
\usepackage[utf8]{inputenc}
\usepackage[russian]{babel}
\usepackage[usenames,dvipsnames]{color}
\usepackage[version=4]{mhchem}
\usepackage[
  a4paper,
  top    = 2.00cm,
  bottom = 2.00cm,
  left   = 1.70cm,
  right  = 1.70cm
]{geometry}
\usepackage{fancyhdr}
\usepackage{hyperref}

\hypersetup{
    colorlinks=true,
    linkcolor=RoyalBlue,
    urlcolor=RoyalBlue,
    citecolor=RoyalBlue,
    anchorcolor=RoyalBlue
}

\urlstyle{same}
\input{glyphtounicode}

\renewcommand{\baselinestretch}{1.5}
\renewcommand{\headrulewidth}{0pt}
\renewcommand{\footrulewidth}{0pt}

\setlength{\footskip}{15pt}
\setlength{\skip\footins}{15pt}
\pagestyle{fancy}

\fancyhf{}
\fancyfoot[R]{\thepage}
\begin{document}

\begin{center}
    {\Huge {\fontsize{25}{30}\selectfont{Реза}} {\fontsize{25}{30}\selectfont{\textbf{Алиасгари Ренани}}}} \\ 
    \vspace{3pt}
    {\small \raisebox{-0.2\height}\ {01/12/2025}} $|$
    {\small \raisebox{-0.2\height}\ {Российская Федерация}} $|$
    {\small \raisebox{-0.2\height}\ {Проект исследования}} $|$
    {\small \raisebox{-0.2\height}\ {Аспирантура}}
    \\
    {\small \raisebox{-0.2\height}\ {Системный анализ, управление и обработка информации, статистика}} 
    \vspace{-10pt}
\end{center}
\vspace{-10pt}

\noindent\rule{\textwidth}{1pt}

Моя исследовательская цель — продвинуть радиационно-стойкие видеопроцессоры на базе FPGA для космических применений путём изучения эффектов ионизирующей радиации на компонентах FPGA и исследования интеграции радиационно-стойкой памяти, такой как ReRAM. Это будет включать детальный анализ эффектов одиночных событий (SEEs) и общего ионизирующего дозового воздействия (TID), за которым последует разработка оптимизированных конвейеров DSP для видео, адаптированных для суровых космических сред. Программа аспирантуры по специальности «Системный анализ, управление и обработка информации, статистика» в МФТИ предоставляет ресурсы и экспертизу для реализации этого, опираясь на мою предыдущую работу в характеризации полупроводников и дизайне FPGA.

\vspace{10pt}

Ионизирующая радиация в космосе, от источников таких как космические лучи и солнечные частицы, представляет значительные вызовы для надёжности FPGA. Ключевые эффекты включают одиночные события переворотов (SEUs) в конфигурационной памяти, которые могут перевернуть биты в ячейках SRAM, приводя к функциональным ошибкам; одиночные события защёлкивания (SELs), которые могут вызвать высокие токи и потенциальный отказ устройства; и накопление общего ионизирующего дозового воздействия, которое со временем ухудшает производительность транзисторов через накопление заряда в оксидах. Наиболее уязвимые компоненты включают конфигурационную логику на базе SRAM, встроенные блоки RAM и MOSFET, где масштабирование к меньшим узлам увеличивает уязвимость к этим эффектам. Мой план начинается с всестороннего тестирования для картирования этих воздействий: использование ускорителей частиц для симуляции космической радиации, измерение скоростей переворотов в различных блоках FPGA и применение техник, таких как профилирование ёмкость-напряжение для оценки плотности ловушек в оксидах после иррадиации.

\vspace{10pt}

Опираясь на это, я намерен исследовать стратегии смягчения, в частности, замену уязвимого SRAM на ReRAM (резистивную RAM) для конфигурационной памяти. Неволатильная природа ReRAM и более высокая стойкость к радиации — благодаря зависимости от изменений сопротивления, а не от хранения заряда — делает её перспективной. Исследования показывают, что интеграция ReRAM поверх слоёв CMOS может уменьшить площадь FPGA в 2-3 раза, одновременно снижая потребление энергии примерно на 20\%, поскольку устраняет необходимость в постоянном обновлении и предлагает лучшую стойкость к SEUs. Я буду исследовать гибридные архитектуры, где ячейки ReRAM хранят конфигурационные данные, потенциально используя 3D-стэкинг для поддержания плотности. Это включает моделирование взаимодействий радиации с материалами ReRAM, такими как оксид гафния, для количественной оценки толерантности к TID и сечений SEU, затем прототипирование маломасштабных тканей FPGA с переключателями ReRAM через инструменты симуляции, такие как Vivado и Simulink.

\vspace{10pt}

Для других компонентов я оценю интеграцию радиационно-стойких альтернатив, таких как использование ячеек памяти большего размера или кодов коррекции ошибок (ECC) в BRAM, и применение тройной модульной избыточности (TMR) к логическим путям, наиболее подверженным транзиентам. Это согласуется с радиационно-стойкими дизайнами FPGA от поставщиков, таких как Xilinx, где улучшенная чистка обнаруживает и исправляет ошибки конфигурации в реальном времени. Мои эксперименты будут включать подвергание модифицированных прототипов FPGA гамма-лучам и тяжёлым ионам, мониторинг функциональной целостности с использованием встроенных самотестов и сравнение метрик производительности, таких как скорости переворотов и потребление энергии, с стандартными устройствами на базе SRAM.

\vspace{10pt}

Заключительная фаза сосредоточена на разработке алгоритмов DSP видео-конвейера, оптимизированных для этих радиационно-стойких FPGA, ориентированных на космические применения, такие как спутниковая съёмка или системы зрения роверов. Опираясь на мой опыт с алгоритмами, такими как детекция краёв Собеля, демозаика и коррекция рыбьего глаза, я расширю на специфические для космоса нужды: снижение шума от артефактов космических лучей, сжатие в реальном времени для передач с ограниченной пропускной способностью и адаптивную фильтрацию при изменяющемся освещении. Эти будут реализованы в конвейерном Verilog, используя срезы DSP FPGA для высокоскоростных операций, с арифметикой фиксированной точки для минимизации использования ресурсов. Для обеспечения толерантности к радиации конвейеры будут включать fault-tolerant дизайны, такие как схемы голосования в установках TMR для маскировки SEUs, и динамическую реконфигурацию для обхода повреждённых блоков.

\vspace{10pt}

Верификация будет включать тестирование hardware-in-the-loop: \\ синтез дизайнов на радиационно-стойких платах, потоковую передачу видео от эмулированных космических камер и подвергание их симулированной радиации через инструменты инъекции ошибок. Цели производительности включают достижение как минимум 60 кадров в секунду при разрешении 1080p с задержкой менее 10 мс, сохраняя функциональность до 100 крад TID. Эта работа завершится интеграцией технологии в российскую промышленность и академию, потенциально сотрудничая с организациями, такими как Роскосмос, для развертывания в орбитальных миссиях, способствуя продвижению в автономной навигации и дистанционном зондировании. Решая эти вызовы, моё исследование повысит надёжность FPGA для аэрокосмической отрасли, поддерживая космические инициативы России через практические, развертываемые системы, и стремится произвести публикации по гибридам ReRAM-FPGA и видеопроцессорам на базе FPGA космического класса.


\end{document}